We review the inference problem and existing SB work that will be relevant to our method. Though our method is more general, we continue the example above to motivate the setup; that is, suppose we are interested in sampling the trajectories of mRNA concentration in a population of cells.

\textbf{Data.}  We have observations at $\totsteps$ time points. At the $\timeidx$th time, $\timestep{\timeidx}$, we observe data from $\tottrajec_{\timestep{\timeidx}}(> 0)$ cells. The mRNA concentration of the $\sampleidx_{\timeidx}$th cell at $\timestep{\timeidx}$ is $\obsat{\timestep{\timeidx}}^{\sampleidx_{\timeidx}}\in \R^d$. After gathering data for a cell, the cell dies, so each cell is observed only once. More generally, we say that we observe the $(\timeidx, \sampleidx_{\timeidx})$ \emph{particle} just once, at time $\timestep{\timeidx}$. Without loss of generality, we choose the time steps to be unique, increasing, and starting at 0: $0 =\timestep{1} < \timestep{2} < \dots < \timestep{\totsteps} < \infty$. They need not be equally spaced.
We let $\obsatall{\timestep{\timeidx}}$ denote all $\tottrajec_{\timestep{\timeidx}}$ observations at time $\timestep{\timeidx}$. There are $\totalnumparticles = \sum_{\timeidx=1}^\totsteps \tottrajec_{\timestep{\timeidx}}$ total observations.

\textbf{Goal.} 
We assume that, if not measured, each cell's mRNA concentration would have had a continuous trajectory in time. $\responseat{\timet}^{(\timeidx, \sampleidx_{\timeidx}) }$ denotes the trajectory for the $\sampleidx_{\timeidx}$th cell observed at the $\timeidx$th time step. So $\obsat{\timestep{\timeidx}}^{\sampleidx_{\timeidx}} = \responseat{\timet = \timestep{\timeidx}}^{ (\timeidx, \sampleidx_{\timeidx}) }.$ That is, $\obsat{\timestep{\timeidx}}^{\sampleidx_{\timeidx}}$ denotes an evaluation of $\responseat{\timet}^{ (\timeidx, \sampleidx_{\timeidx}) }$ at the observed time $\timestep{\timeidx}$. We assume the trajectories are independent samples from a latent distribution over trajectories; this assumption implies the observations are independent as well. Our goal is to generate samples from the latent distribution over particle trajectories within $\timet \in [0,\timestep{\totsteps}]$, given the observations $\{ \obsat{\timestep{\timeidx}}^{\sampleidx_{\timeidx}} \}_{\timeidx, \sampleidx_{\timeidx}}$.

\textbf{Model.}
We model the latent trajectory of the $(\timeidx, \sampleidx_{\timeidx})$ particle with an SDE driven by a $d$-dimensional Brownian motion $\brownianm^{ (\timeidx,\sampleidx_{\timeidx}) }$, independent across particles:
\begin{equation}
    \de\responseat{\timet}^{(\timeidx, \sampleidx_{\timeidx}) } = \modeldrift(\responseat{\timet}^{ (\timeidx,\sampleidx_{\timeidx}) },t) \de\timet + \sqrt\volatility \de\brownianm^{ (\timeidx,\sampleidx_{\timeidx}) },~~\responseat{\timet = 0}^{ (\timeidx,\sampleidx_{\timeidx}) }\sim \marginal{0}.
    \label{eq:mainsde}
\end{equation}
We assume that the volatility $\volatility$ is known.\footnote{This assumption is common in the SB literature since it ensures that the SB problem is well-posed \citep{Chen2022, Vargas2021, lavenant2024toward}. In our experiments, we use the same fixed volatility value as past SB work, and we find it works well. Estimating volatility from data is an interesting direction for future research.} We assume that the drift $\modeldrift(\cdot,\cdot): \R^{d}\times [0,\timet_{\totsteps}] \to \R^{d}$ and initial marginal distribution $\marginal{0}$ are unknown.

We assume standard SDE regularity conditions. The first assumption below ensures a strong solution to the SDE exists; see \citet[][Chapter 3, Theorem 3.1]{pavliotis2016stochastic}. The second ensures that the process does not exhibit unbounded variability.
\begin{assumption}
\label{assumption-lipschitz}
The drifts are $L$-Lipschitz; i.e., for all $t\in [0,\timet_{\totsteps}]$, $\|\modeldrift(x, t)-\modeldrift(y, t)\| \le L \| x-y \|$, where $\|\cdot \|$ denotes the usual $L^2$ norm of a vector. And we have at most linear growth; i.e., there exist $K<\infty$ and constant $c$ such that $\| \modeldrift(y, t)\| <K \| y \|+c$. 
\end{assumption} 
\begin{assumption}
\label{assumption-bddsecondmoments}
At each time step $\timestep{\timeidx}$, the distribution of the $\tottrajec_{\timestep{\timeidx}}$ particles has bounded second moments.
\end{assumption}

\textbf{Generating sample trajectories.} One approach to generating trajectory samples is to first estimate the unknown drift from data. Using the estimated drift, we can simulate forward and backward in time starting from a particle observation to generate an approximate realization of the SDE solution for $\timet \in [0,\timestep{\totsteps}]$. This approach generates one sample trajectory for each observation, for a total of $\totalnumparticles$ trajectory samples.

More precisely, the drift $\modeldrift$ in \cref{eq:mainsde} defines the forward dynamics of $\responseat{\timet}^{ (\timeidx,\sampleidx_{\timeidx}) }$. Given \cref{eq:mainsde}, the backwards process $\responseat{\timet_{\totsteps}-\timet}^{ (\timeidx,\sampleidx_{\timeidx}) }$ is described by an analogous SDE with the same volatility and a new (backward) drift, $\modeldriftbackward(\cdot,\cdot)$; see \citet{haussmann1986time, follmer2005entropy, cattiaux2023time} for details. Suppose we had access to estimates $\estimatedrift$ and $\estimatebkwdrift$ for the forward and backward drift, respectively. Then given a particle observed at time $\timestep{\timeidx}$, we could simulate a trajectory forward in $\timet \in [\timestep{\timeidx},\timestep{\totsteps}]$ by using \cref{eq:mainsde} with (1) $\modeldrift=\estimatedrift$, (2) the initial distribution equal to a point mass at the observation, and (3) discretized time steps of size $\Delta\timet$. We write the resulting simulated trajectory as $\simulatedresponseat{\timestep{\timeidx}\le \timet\le \timestep{\totsteps}}^{ (i,\sampleidx_{\timeidx}) }=\textrm{\texttt{forwardSDE}}(\estimatedrift, \obsat{\timestep{\timeidx}}^{\sampleidx_{\timeidx}},\Delta \timet)$. 
Similarly, we can sample the part of the trajectory within $\timet \in [0,\timestep{\timeidx}]$ using the backward drift and SDE: $\simulatedresponseat{0\le \timet\le \timestep{\timeidx}}^{(i,\sampleidx_{\timeidx})}=\textrm{\texttt{backwardSDE}}(\estimatebkwdrift, \obsat{\timestep{\timeidx}}^{\sampleidx_{\timeidx}},\Delta \timet)$. In our experiments, we use a standard Euler-Maruyama approach \citep[][Chapter 3.4]{sarkka2019applied}; for full details, see \cref{sec:sampling-routine}.  It remains to estimate the forward and backward drifts.

\textbf{Multi-marginal Schrödinger bridges.}
When we have some prior knowledge about the system's latent dynamics, Schrödinger Bridges \citep[SBs,][]{Wang2023, Vargas2021, lavenant2024toward} can be used to estimate the forward and backward drift. We next review multi-marginal SBs and how their drift estimates arise. The general multi-marginal SB problem \citep{chen2019multi, lavenant2024toward} finds the \emph{bridge} distribution $\bridge$ over trajectories that (1) interpolates between specified marginals while (2) being closest --- in a Kullback--Leibler divergence ($\KL$) sense --- to the prior knowledge as expressed in a \emph{reference} distribution $\refden$. With a slight abuse of notation, we use $\bridge$ and $\refden$ to denote not just the distributions over trajectories but also their respective densities with respect to a Wiener measure. Let the distribution over trajectory points at time $\timestep{\timeidx}$ implied by $\bridge$ be $\popmarginalat{\timestep{\timeidx}}$. Let $\marginal{\timestep{\timeidx}}$ be a desired marginal at time $\timestep{\timeidx}$. Let $[I] := \{1,\ldots, I\}$. Then the multi-marginal SB problem is a constrained optimization problem:\footnote{In the SB literature, this KL divergence is often defined in terms of probability distributions rather than densities. But implicitly, to compute the KL divergence, one uses the densities. So here we directly use the densities.} 
\begin{align}
    \label{eq:SBP-main}
    & \argmin_{\bridge\in \measurefamily: \bridge_0=\refden_0}\KL(\bridge||\refden), \\
    \label{eq:measure-family}
    \textrm{where} & \;\; \measurefamily=\left\{\bridge:\forall i \in [I], \popmarginalat{\timestep{\timeidx}}=\marginal{\timestep{\timeidx}} \right\}
\end{align}
is the family of densities over trajectories satisfying the marginal constraints.

Suppose now that the reference is a distribution over trajectories implied by an SDE as in \cref{eq:mainsde}. Since we assume the volatility $\volatility$ is known, we assume the known volatility is used in the reference. So the reference is specified by its drift $\refdrift$. Under \cref{assumption-lipschitz,assumption-bddsecondmoments}, the resulting probability distribution over trajectories admits a density with respect to the Wiener measure defined by a Brownian motion with volatility $\volatility$ \citep{oksendal2013stochastic, kailath1971structure}. The existence of the density follows from Girsanov's Theorem; see, e.g., Theorem 8.6.3 in \citet{oksendal2013stochastic}) and \cref{app:convex-densities}. We write $\refden = \refden_{\refdrift}$ for the distribution and its density. In this case, \citet{kailath1971structure} shows that not only must the optimal $\bridge$ also solve an SDE of the form in \cref{eq:mainsde}, but it must have the same volatility; $\KL$ will be infinite for other volatilities. Therefore, the optimal $\bridge$ is defined by its drift, and we write $\bridge = \bridge_{\modeldrift}$. Then solving \cref{eq:SBP-main} and \cref{eq:measure-family} is equivalent to finding the best drift for $\bridge_{\modeldrift}$:
\begin{align}
    \label{eq:SBP-drift}
    \argmin_{\modeldrift: \bridge_{\modeldrift} \in \measurefamily,  \bridge_{\modeldrift,0}=\refden_{\refdrift,0} }\KL(\bridge_{\modeldrift}||\refden_{\refdrift} ),
\end{align}
with $\measurefamily$ as in \cref{eq:measure-family}. Practical SB algorithms take the marginals to be the (discrete) empirical distributions of observed data at each time point; that is, $\marginal{\timestep{\timeidx}}$ is constructed from $\obsatall{\timestep{\timeidx}}$, where  $\obsatall{\timestep{\timeidx}}$ represents the collection of all the observations available at time $\timestep{\timeidx}$ \citep{Vargas2021, de2021diffusion}.

Even though \cref{eq:SBP-drift}  appears to leverage information from all time snapshots jointly, \citet{lavenant2024toward, chen2019multi} showed that --- under very mild conditions --- this SB problem is equivalent to a collection of $\totsteps-1$ separate SB problems, each between two adjacent time points. Therefore, practical algorithms for solving this problem return $\totsteps-1$ pairs of forward and backward drifts $\{\estimatedriftstep{i},\estimatebackwarddriftstep{i}\}_{\timeidx=1}^{\totsteps-1}$, where $\estimatedriftstep{i}$ and $\estimatebackwarddriftstep{i}$ are defined over the interval $[\timestep{\timeidx}, \timestep{\timeidx+1}]$. We can construct an estimate $\estimatedrift$ of the forward drift in the time span $\timet \in [0,\timestep{\totsteps}]$ by concatenating these estimates, with an analogous estimate $\estimatebkwdrift$ for the backward drift.

\textbf{Challenges.} Since this approach performs piecewise interpolation between consecutive pairs 
it does not share information across different time intervals. Moreover, this approach requires  practitioners to specify a single reference drift $\refdrift$. Even when scientists know a parametric form for the underlying dynamical system \citep[e.g.,][for mRNA dynamics]{pratapa2020benchmarking}, they rarely have access to all parameter values. So in practice the choice of $\refdrift$ usually corresponds to Brownian motion \citep{Caluya2022, Vargas2021,Chen2022, chen2019multi, chen2024deep}.





